\chapter{Bond and ETF Request-for-Quote Market Making}\label{ch:hf:bond-and-etf-request-for-quote-market-making}

A corporate bond that last traded three weeks ago is on the screen, in a request from an asset manager to five dealers. One of them answers in
under a second, without a human, because its model has already priced the bond from its neighbours. In this chapter's auction, five dealers whose
estimates of such a bond are each off by a quarter of a point settle on markups of about 43 cents; the winner makes 0.94 cents per request it
answers and wins 13\% of them. A dealer that sets its markup as if its estimate were right would win 31\% and lose money.

\section{Pricing bonds that do not trade}

Most corporate bonds trade rarely (One Quant Book 2, chapter 22): an issue may print a few times a month, a request for quote may be the only price
discovery it gets in a week. A dealer's fair value for such a bond is assembled, not observed. Three sources feed it:

\begin{itemize}
\item \emph{its own prints}, from the trade reporting system, weighted by recency and cleaned of outliers into a composite (Book 2's
\texttt{firm.rfq.composite}), and moved since by the issuer's curve;
\item \emph{comparables}: the same issuer's other bonds, bonds of similar rating, sector and maturity, each with its sensitivity to the target;
\item \emph{liquid proxies}: the Treasury curve, credit indices and bond exchange-traded funds, which move every second and carry the market-wide
part of the bond's value.
\end{itemize}

\texttt{firm.rfqmm.fair\_value} blends the first two. Whatever the method, the estimate has an error, and for a bond that has not traded for weeks
the error is large relative to the bid--ask spread a dealer can charge. The chapter's model takes it at 25 cents per 100 of face value, one standard
deviation, for every dealer; each dealer's error is its own.

\section{Auto-quoting request for quote}

\begin{definition}[Auto-quoting]\label{def:hf:bond-and-etf-request-for-quote-market-making:auto}
\emph{Auto-quoting}\index{auto-quoting} is a dealer's automatic response to requests for quote: a system that prices each request from its models
and inventory, decides whether and at what level to answer, and sends the answer without a trader, within limits set by the desk and with requests
outside them routed to a person.
\end{definition}

\begin{omfigure}
\begin{tikzpicture}[font=\footnotesize,>=stealth,box/.style={draw,rounded corners,fill=omThm!8,minimum width=2.2cm,minimum height=0.7cm,align=center}]
\node[box] (c) at (0,0) {client\\(sells \$5m)};
\node[box] (p) at (3.2,0) {platform};
\foreach \i/\y in {1/1.6,2/0.8,3/0,4/-0.8,5/-1.6}{
  \node[box,minimum width=1.6cm] (d\i) at (6.8,\y) {dealer \i};
  \draw[->] (p.east) -- (d\i.west);
}
\node[box,fill=omDef!10] (m) at (10.3,0) {fair value\\$-$ markup\\(auto-quote)};
\draw[->] (d3.east) -- (m.west);
\draw[->] (c) -- node[above] {request} (p);
\draw[->,dashed] (d1.north) to[bend right=12] (c.north);
\node[font=\scriptsize,fill=white,inner sep=1pt] at (3.4,2.25) {best bid wins; the winner sees the cover};
\end{tikzpicture}

\omcaption{A request for quote on a dealer-to-client platform: the client asks five dealers at once; each prices the bond from its own model and
answers within the time limit; the client hits the best bid if it is acceptable; the platform tells the winner the cover, the second-best bid.
Source: Book 2, chapter 22; this chapter.}\label{fig:hf:bond-and-etf-request-for-quote-market-making:flow}
\end{omfigure}

Hendershott and Madhavan studied the innovation at the heart of it (\cref{fig:hf:bond-and-etf-request-for-quote-market-making:flow}): electronic
technology that lets an investor search many bond dealers at once. They showed that periodic one-sided electronic auctions are a viable and
important source of liquidity even in inactively traded bonds, a compromise between bilateral search and a continuous order book. O'Hara and Zhou,
combining a platform's proprietary data with the regulatory trade reports, found electronic request-for-quote trading still fairly small and
segmented, but with wide-ranging effects on transaction costs and execution quality in both electronic and voice trading, and in the interdealer
market: a market in transition.

\begin{dated}{2026-09}{One platform's year}\label{dat:hf:bond-and-etf-request-for-quote-market-making:mktx}
MarketAxess reported record revenue of \$846 million for 2025, a 48\% increase in portfolio-trading average daily volume to a record \$1.4 billion,
a 24\% increase in block-trading average daily volume to a record \$5 billion and a 33\% increase in dealer-initiated volume; about 2,100 firms use
its platform.
\end{dated}

\section{Win probability and the winner's curse}

\begin{definition}[Cover price]\label{def:hf:bond-and-etf-request-for-quote-market-making:cover}
The \emph{cover price}\index{cover price} is the second-best quote in a request for quote, which many platforms disclose to the winning dealer after
the trade; it tells the winner how much it could have shaded its quote and still won.
\end{definition}

\begin{definition}[Win-probability model]\label{def:hf:bond-and-etf-request-for-quote-market-making:win}
A \emph{win-probability model}\index{win-probability model} estimates the probability that a dealer's quote wins a request for quote, as a function
of the quote's distance from the dealer's fair value and of the request's features (size, side, client, bond, time since the last print, number of
dealers asked), fitted on the dealer's history of wins, losses and covers.
\end{definition}

The dealer chooses its markup $m$, the distance of its bid below its fair value, to maximise $P(\text{win}\mid m)\times\mathrm{E}[\text{profit}\mid
\text{win},m]$. The second factor is where the winner's curse lives: the dealer whose estimate is highest wins most often, so its estimate, given that
it won, is too high. With five dealers each off by 25 cents, the highest estimate is 29 cents too high on average before any markup (the expected
largest of five normal errors, Book 2's \texttt{firm.rfq.expected\_max\_normal}); with ten, 38.5 cents.

The chapter's auction (\cref{lst:hf:bond-and-etf-request-for-quote-market-making:sim}): the client sells to the best bid if it is at least its
\omterm{def:hf:inventory-models:reservation}{reservation price}, the bond's value less an urgency discount with a mean of 30 cents. Competitors use a common markup; ours is the best response;
the equilibrium is the markup at which the best response to everyone using it is itself. A model that ignores the curse maximises $P(\text{win})\times m$,
as if the profit on a win were the markup.

\begin{omfigure}
\begin{tikzpicture}
\begin{axis}[width=11cm,height=5.4cm,axis y line*=right,axis x line=none,xmin=0.5,xmax=10.5,ymin=0,ymax=32,
  ylabel={\small win rate (\%)},tick label style={font=\footnotesize},table/col sep=comma]
\addplot[gray,thick,dotted,mark=triangle*,mark size=1.4pt] table[x=n,y=win]{figdata/hft/22-bond-and-etf-request-for-quote-market-making/dealers.csv};
\label{plot:hf:rfqwin}
\end{axis}
\begin{axis}[width=11cm,height=5.4cm,axis y line*=left,xmin=0.5,xmax=10.5,xtick={1,2,3,5,8,10},xlabel={\small dealers asked},
  ylabel={\small cents per request (per 100 face)},ymin=-2,ymax=9,tick label style={font=\footnotesize},grid=major,grid style={gray!25},
  table/col sep=comma,legend style={font=\scriptsize,at={(0.5,-0.3)},anchor=north},legend columns=3]
\addplot[omThm,thick,mark=*,mark size=1.3pt] table[x=n,y=pnl]{figdata/hft/22-bond-and-etf-request-for-quote-market-making/dealers.csv};
\addlegendentry{profit at the equilibrium markup}
\addplot[omDef,thick,dashed,mark=square*,mark size=1.3pt] table[x=n,y=naive]{figdata/hft/22-bond-and-etf-request-for-quote-market-making/dealers.csv};
\addlegendentry{ignoring the curse}
\addlegendimage{/pgfplots/refstyle=plot:hf:rfqwin}\addlegendentry{win rate (right)}
\draw[gray] (axis cs:0.5,0) -- (axis cs:10.5,0);
\end{axis}
\end{tikzpicture}

\omcaption{A dealer's expected profit per request answered, in cents per 100 of face value, at the symmetric equilibrium markup and at the markup a
model ignoring the winner's curse would choose (competitors at the equilibrium), and the equilibrium win rate (right), against the number of dealers
asked; estimates off by 25 cents, urgency discounts with a mean of 30 cents, 40,000 requests per point. Data:
\texttt{hf\_rfq.equilibria}.}\label{fig:hf:bond-and-etf-request-for-quote-market-making:dealers}
\end{omfigure}

\begin{center}
\small
\begin{tabular}{@{}lrrrrrr@{}}
\toprule
dealers asked & 1 & 2 & 3 & 5 & 8 & 10\\
\midrule
equilibrium markup (cents) & 47.5 & 40.0 & 40.0 & 42.5 & 45.0 & 50.0\\
profit per request (cents) & 7.96 & 3.46 & 1.96 & 0.94 & 0.41 & 0.45\\
win rate & 27.7\% & 24.0\% & 19.1\% & 13.2\% & 9.1\% & 6.9\%\\
winner's error given a win (cents) & 18.8 & 25.6 & 29.7 & 35.4 & 40.5 & 43.5\\
markup ignoring the curse (cents) & 37.5 & 30.0 & 27.5 & 25.0 & 22.5 & 25.0\\
its profit per request (cents) & 7.50 & 2.99 & 1.25 & $-0.07$ & $-1.13$ & $-0.65$\\
\bottomrule
\end{tabular}
\end{center}

The results (\cref{fig:hf:bond-and-etf-request-for-quote-market-making:dealers}): the equilibrium markup does not fall with more dealers, it rises
after three, because the curse grows faster than competition narrows the spread; profit per request falls from 8 cents alone to under half a cent
with eight or ten. The dealer that ignores the curse quotes 20 to 25 cents tighter, wins more than twice as often, and loses from five dealers on.
Even alone it earns less, because a single dealer's winning bids are also selected, by the client's \omterm{def:hf:inventory-models:reservation}{reservation price}.

The \omterm{def:hf:bond-and-etf-request-for-quote-market-making:win}{win-probability model} is learned from history. Fitting a logistic curve to 20,000 past requests with markups spread from 0 to 80 cents against
five competitors at the equilibrium (\cref{lst:hf:bond-and-etf-request-for-quote-market-making:win}) gives
$P(\text{win})=1/(1+e^{-(0.665-0.0612\,m)})$: 12.6\% at the equilibrium markup, against 13.2\% simulated
(\cref{fig:hf:bond-and-etf-request-for-quote-market-making:winmodel}). The cover feeds the model: each win reveals how far the dealer could have
shaded, each loss nothing, which biases naive estimates of the distribution of competitors' quotes; models of it are fitted with that censoring in
mind.

\begin{omfigure}
\begin{tikzpicture}
\begin{axis}[width=11cm,height=4.8cm,xlabel={\small markup below fair value (cents per 100 face)},ylabel={\small probability of winning},xmin=0,
  xmax=80,ymin=0,ymax=0.65,tick label style={font=\footnotesize},grid=major,grid style={gray!25},table/col sep=comma,
  legend style={font=\scriptsize,at={(0.97,0.97)},anchor=north east}]
\addplot[omThm,only marks,mark=*,mark size=1.6pt] table[x=markup,y=empirical]{figdata/hft/22-bond-and-etf-request-for-quote-market-making/winmodel.csv};
\addplot[omDef,thick] table[x=markup,y=fitted]{figdata/hft/22-bond-and-etf-request-for-quote-market-making/winmodel.csv};
\legend{win rate in 10-cent bins,fitted logistic model}
\end{axis}
\end{tikzpicture}

\omcaption{The \omterm{def:hf:bond-and-etf-request-for-quote-market-making:win}{win-probability model}: win rates of 20,000 historical requests grouped by markup, and the logistic curve fitted to them by Newton's
method, against five dealers at the equilibrium markup of 42.5 cents. Data: \texttt{hf\_rfq.win\_model}.}\label{fig:hf:bond-and-etf-request-for-quote-market-making:winmodel}
\end{omfigure}

\section{ETF request for quote and block liquidity}

Institutions trade large blocks of exchange-traded funds by request for quote, off the exchange's order book. The dealer's bid for a block comes
from the creation basket (chapter 12): it can buy the shares and redeem them, receiving the basket, which it sells. Its bid is the basket's bid less
the redemption's costs: the fixed fee spread over the unit's shares, the basket's selling cost and the hedge while the redemption settles. For a fund
whose basket bids \$100.00, with units of 50,000 shares, a \$1,250 fee, 5 basis points to sell the basket and 1 to hedge, the dealer can bid
\$99.915 a share for a block of any size that fits its balance sheet. The winner's curse is small here, because the basket is priced continuously;
the risk is in the basket's liquidity when the block is large.

\section{Portfolio trades and all-to-all}

A portfolio trade (Book 2, chapter 22) prices a list of hundreds of bonds as one package, often against an exchange-traded fund that holds most of
them: the dealer's model prices every line, and the package's diversification shrinks the winner's curse on the whole relative to the sum of its
parts. The dated box's growth in portfolio trading is the market choosing that trade. All-to-all trading (Book 2, chapter 22) lets investors answer
each other's requests anonymously: non-dealers become liquidity providers, and the auto-quoter competes with them.

\section{Strategy files}

\begin{strategyfile}{Algorithmic corporate-bond request-for-quote responder}\label{strat:hf:bond-and-etf-request-for-quote-market-making:corp}
\sfield{Who pays you, and why} Investors who need to trade bonds that rarely trade, and pay for a dealer's balance sheet and estimate.
\sfield{Instruments and venues} Corporate bonds on dealer-to-client platforms.
\sfield{Signal} A fair value from the bond's stale prints, comparables and liquid proxies; the \omterm{def:hf:bond-and-etf-request-for-quote-market-making:win}{win-probability model}.
\sfield{Sizing and execution} Answer automatically within limits; set the markup from the model with the winner's curse; route large or unusual
requests to a trader.
\sfield{Costs} The curse (35 cents of error given a win against five dealers in the model); inventory financing and hedging.
\sfield{How it dies} Markups set as if the estimate were right (with five dealers the naive quoter wins 31\% of requests and loses money); and
stress, when every estimate goes stale at once, as in the corporate bond dislocation of March 2020 (One Quant Book 9, chapter 19).
\sfield{Horizon, capacity, infrastructure} Seconds to answer, days to weeks to unwind; the platform's connections and trade reports.
\sfield{Backtest honestly} The requests the dealer actually received, the quotes competitors made (censored by what was disclosed), the prices at
which the inventory could be sold.
\sfield{Sources} Hendershott and Madhavan (2015); O'Hara and Zhou (2021); this chapter.
\end{strategyfile}

\begin{strategyfile}{ETF request-for-quote market making}\label{strat:hf:bond-and-etf-request-for-quote-market-making:etf}
\sfield{Who pays you, and why} Institutions trading fund blocks larger than the exchange's displayed depth.
\sfield{Instruments and venues} Exchange-traded funds by request for quote; their baskets; the primary market.
\sfield{Signal} The basket's executable value and the costs of creating or redeeming.
\sfield{Sizing and execution} Quote from the basket less (or plus) the primary market's costs; hedge in the basket or a future at once.
\sfield{Costs} \omterm{def:hf:etf-market-making:fee}{Creation fees}, the basket's spread, the hedge.
\sfield{How it dies} Illiquid baskets in a stress, where the basket's bid is itself an estimate: the bond-fund discounts of March 2020 (One Quant
Book 9, chapter 19) and chapter 12's simulated freeze.
\sfield{Horizon, capacity, infrastructure} Seconds to a day; the fund's creation cut-off.
\sfield{Backtest honestly} Basket prices at the request's time; the fee schedule; settlement.
\sfield{Sources} Chapter 12; the dated box.
\end{strategyfile}

\begin{strategyfile}{Treasury streaming and request for quote}\label{strat:hf:bond-and-etf-request-for-quote-market-making:ust}
\sfield{Who pays you, and why} Investors trading Treasuries on dealer-to-client platforms, who prefer a disclosed dealer's firm price.
\sfield{Instruments and venues} On-the-run and off-the-run Treasuries; streams and requests on dealer-to-client platforms.
\sfield{Signal} The interdealer order books (One Quant Book 2, chapter 4) for on-the-runs; a fitted curve for off-the-runs.
\sfield{Sizing and execution} Stream tight on on-the-runs, hedge in the interdealer market or futures; answer requests on off-the-runs from the curve.
\sfield{Costs} Hedge spreads; curve-fit error on off-the-runs.
\sfield{How it dies} Competitors with faster interdealer access, and days when the interdealer book thins: on 15 October 2014 the ten-year yield
fell 16 basis points between 9:33 and 9:39 a.m. and nearly retraced it by 9:45, with no apparent trigger and uncharacteristically shallow depth
(the official joint staff report of 2015).
\sfield{Horizon, capacity, infrastructure} Milliseconds to hours.
\sfield{Backtest honestly} Interdealer prices at each request's time, not the day's close.
\sfield{Sources} One Quant Book 2, chapter 4; Joint Staff Report, The U.S. Treasury market on October 15, 2014 (2015).
\end{strategyfile}

\begin{strategyfile}{All-to-all anonymous liquidity provision}\label{strat:hf:bond-and-etf-request-for-quote-market-making:a2a}
\sfield{Who pays you, and why} Investors whose requests reach non-dealers anonymously.
\sfield{Instruments and venues} Corporate bonds on all-to-all protocols.
\sfield{Signal} The same fair value, with no client identity to price.
\sfield{Sizing and execution} Answer where the model's estimate is best (bonds with recent prints, close comparables); size small.
\sfield{Costs} A curse sharper than with known clients, since the flow cannot be tiered.
\sfield{How it dies} Adverse selection from better-informed anonymous requesters.
\sfield{Horizon, capacity, infrastructure} Seconds to days.
\sfield{Backtest honestly} Anonymous requests' later price moves by requester type where the platform reports it.
\sfield{Sources} O'Hara and Zhou (2021); Book 2, chapter 22.
\end{strategyfile}

\section{Tutorial: five dealers and a bond nobody traded}\label{tut:hf:bond-and-etf-request-for-quote-market-making:tut}

\begin{tutorial}
\textbf{Goal.} Find the equilibrium markup of an auto-quoter against competitors with the same estimation error, measure the winner's curse, and fit
a \omterm{def:hf:bond-and-etf-request-for-quote-market-making:win}{win-probability model}. \textbf{End state:} the table and the three figures.
\begin{tutsteps}
\item \textbf{The auction}: every dealer's estimate with its own error, the client's \omterm{def:hf:inventory-models:reservation}{reservation price}, the best bid.
\omcode{code/firm/rfqmm/firm_rfqmm.py}{50}{63}{Estimates, markups, the client's choice, and the winner's profit against the true value.}\label{lst:hf:bond-and-etf-request-for-quote-market-making:sim}
\item \textbf{The equilibrium} by best responses on a grid of markups (\texttt{firm.rfqmm.equilibrium}).
\item \textbf{The win model} fitted on a randomised history.
\omcode{code/firm/rfqmm/firm_rfqmm.py}{97}{108}{A logistic win probability in the markup, by Newton's method.}\label{lst:hf:bond-and-etf-request-for-quote-market-making:win}
\item \textbf{The ETF block bid} from the creation basket (\texttt{firm.rfqmm.etf\_rfq\_bid}).
\end{tutsteps}
\textbf{What to change next.} Give dealers different errors (a better model wins more and profits more); add inventory to the markup (axes);
estimate the competitors' distribution from covers with censoring.
\end{tutorial}

\section{Build: the RFQ market maker}\label{bld:hf:bond-and-etf-request-for-quote-market-making:rfqmm}

\begin{build}
\bfield{Purpose} Price illiquid bonds, answer requests for quote at the markup that maximises expected profit net of the winner's curse, and price ETF
blocks from their baskets.
\bfield{Interface} \texttt{fair\_value(prints, ages\_days, comps, comp\_betas, curve\_move, half\_life)}, \texttt{simulate(n\_dealers, our\_markup, n,
seed, sigma\_est, comp\_markup, urgency)}, \texttt{optimise}, \texttt{equilibrium}, \texttt{naive\_markup}, \texttt{fit\_win\_model},
\texttt{etf\_rfq\_bid}. Built on \texttt{firm.rfq} (Book 2).
\bfield{Rules} Cents per 100 of face value; the client sells; profit against the bond's true value.
\bfield{Acceptance tests} \texttt{code/firm/rfqmm/tests/}: the composite-based fair value within the comparables' range; profit equals value less bid on
wins; the winner's average error is positive; higher markups win less; with five dealers profit and win rate are below one dealer's; the naive
markup is lower and earns less; the logistic fit recovers known coefficients; the ETF bid by hand.
\bfield{Stretch} Heterogeneous dealers; inventory and axes; censored estimation from covers.
\end{build}

\begin{omsources}
\item T.~Hendershott, A.~Madhavan, Click or call? Auction versus search in the over-the-counter market, \emph{Journal of Finance} 70(1), 2015,
419--447.
\item M.~O'Hara, X.~A.~Zhou, The electronic evolution of corporate bond dealers, \emph{Journal of Financial Economics} 140(2), 2021, 368--390.
\item MarketAxess Holdings, fourth quarter and full year 2025 results, Form 8-K exhibit 99.1, 6 February 2026.
\item US Department of the Treasury, Federal Reserve, SEC and CFTC staff, Joint Staff Report: The U.S. Treasury market on October 15, 2014, July 2015.
\end{omsources}

\section{Exercises}

\begin{exercise}[$\star$]\label{exo:hf:bond-and-etf-request-for-quote-market-making:1}
A fund's basket bids \$100.00; units are 50,000 shares; the fee is \$1,250; selling the basket costs 5 basis points and hedging 1. What is the block
bid per share?
\end{exercise}

\begin{exercise}[$\star$]\label{exo:hf:bond-and-etf-request-for-quote-market-making:2}
With estimation errors of 25 cents and five dealers, by how much is the highest estimate too high on average?
\end{exercise}

\begin{exercise}[$\star$]\label{exo:hf:bond-and-etf-request-for-quote-market-making:3}
The fitted model is $P=1/(1+e^{-(0.665-0.0612m)})$. What is the probability of winning at a markup of 42.5 cents? At 20?
\end{exercise}

\begin{exercise}[$\star\star$]\label{exo:hf:bond-and-etf-request-for-quote-market-making:4}
Why does the equilibrium markup rise from five dealers to ten?
\end{exercise}

\begin{exercise}[$\star\star$]\label{exo:hf:bond-and-etf-request-for-quote-market-making:5}
Why does disclosing the cover only to the winner bias a \omterm{def:hf:bond-and-etf-request-for-quote-market-making:win}{win-probability model} fitted naively?
\end{exercise}

\begin{exercise}[$\star\star$]\label{exo:hf:bond-and-etf-request-for-quote-market-making:6}
Why is the winner's curse smaller in a portfolio trade than in the sum of its lines?
\end{exercise}

\begin{exercise}[$\star\star\star$]\label{exo:hf:bond-and-etf-request-for-quote-market-making:7}
\emph{Coding.} With \texttt{firm.rfqmm.optimise}, find our best markup against five competitors at 42.5 cents when our estimate's error is 15
cents and theirs 25. What do we earn and win?
\end{exercise}

\begin{exercise}[$\star\star\star$]\label{exo:hf:bond-and-etf-request-for-quote-market-making:8}
\emph{Find the flaw.} ``Our auto-quoter wins 30\% of requests with a 25-cent markup, and our average markup on wins is 25 cents: it earns 25 cents
a win.''
\end{exercise}

\section{Problem: Five Dealers and a Bond Nobody Traded}

\begin{problem}[{Weekend problem --- five dealers and a bond nobody traded}]\label{pb:hf:bond-and-etf-request-for-quote-market-making:1}
A dealer automates its answers to requests for quote in illiquid corporate bonds.

\medskip\noindent\textbf{Part I --- Prices.}
\begin{enumerate}
\item How is a fair value assembled for a bond that has not traded for weeks?
\item Define \omterm{def:hf:bond-and-etf-request-for-quote-market-making:auto}{auto-quoting}.
\item What did Hendershott and Madhavan and O'Hara and Zhou find?
\item Summarise the dated box.
\end{enumerate}

\medskip\noindent\textbf{Part II --- The auction.}
\begin{enumerate}[resume]
\item Define the \omterm{def:hf:bond-and-etf-request-for-quote-market-making:cover}{cover price} and a \omterm{def:hf:bond-and-etf-request-for-quote-market-making:win}{win-probability model}.
\item Write the dealer's objective and explain the winner's curse in it.
\item How large is the curse for five and ten dealers before markups?
\item How is the equilibrium found?
\end{enumerate}

\medskip\noindent\textbf{Part III --- Measurements.}
\begin{enumerate}[resume]
\item Give the equilibrium markup, profit and win rate for 1, 5 and 10 dealers.
\item What does ignoring the curse do?
\item How well does the logistic model fit?
\item Price the ETF block of the example.
\end{enumerate}

\medskip\noindent\textbf{Part IV --- The verdict.}
\begin{enumerate}[resume]
\item State the \emph{named result}: the auto-quoter's profit per request and win rate at its optimal markup, and how both change with the number of
dealers asked.
\item How many dealers should a client ask?
\item What does a better estimate than the competitors' buy?
\item Why do portfolio trades grow?
\item What changes in all-to-all trading?
\item Which strategy file is most exposed to the curse?
\item How would you estimate competitors' quotes from covers?
\item In one sentence: what does a bond auto-quoter sell?
\end{enumerate}
\end{problem}

\section{Interview questions}

\begin{interviewq}[$\star$ \iqroles{trader}]\label{iq:hf:bond-and-etf-request-for-quote-market-making:1}
You won a request at 99.50 and the cover was 99.20. What does that tell you?
\end{interviewq}

\begin{interviewq}[$\star\star$ \iqroles{researcher}]\label{iq:hf:bond-and-etf-request-for-quote-market-making:2}
How would you price a bond that has not traded in a month?
\end{interviewq}

\begin{interviewq}[$\star\star$ \iqroles{developer}]\label{iq:hf:bond-and-etf-request-for-quote-market-making:3}
Design an auto-quoter that answers 50,000 requests a day within 500 milliseconds each. Where are the risks?
\end{interviewq}

\begin{interviewq}[$\star\star$ \iqroles{risk}]\label{iq:hf:bond-and-etf-request-for-quote-market-making:4}
The auto-quoter's win rate doubled last month. What do you check?
\end{interviewq}

\begin{interviewq}[$\star\star$ \iqroles{trader}]\label{iq:hf:bond-and-etf-request-for-quote-market-making:5}
A client asks you and nineteen other dealers. How do you quote?
\end{interviewq}

\begin{interviewq}[$\star\star\star$ \iqroles{researcher}]\label{iq:hf:bond-and-etf-request-for-quote-market-making:6}
With $n$ dealers whose estimates have independent normal errors of standard deviation $\sigma$, derive the expected error of the highest estimate for
$n=2$.
\end{interviewq}
